% ============================================================================
%  CAPÍTULO III: DESARROLLO DEL TRABAJO (parte 3)
% ============================================================================

\section{Consideraciones del sistema operativo destino}

\subsection{Aislamiento con contenedores}

El entorno se ejecuta sobre un host Linux con Docker Engine y Docker Compose
v2. Cada servicio corre en su propio contenedor, con su propio sistema de
archivos, su propia pila de red y sus propios límites de recursos. Los
contenedores comparten el núcleo del host, pero no comparten el espacio de
usuarios ni los procesos. Esa separación reduce el impacto de una
vulnerabilidad en un servicio y facilita la reproducibilidad
\citep{docker2024}.

Los puertos se publican únicamente en \texttt{127.0.0.1}. Esa decisión evita
que MySQL o PostgreSQL queden accesibles desde la red local o desde Internet.
Si un integrante necesita conectarse con una herramienta gráfica, puede usar
un túnel SSH o una regla explícita; no se abre el puerto a \texttt{0.0.0.0}
por comodidad. La red interna \texttt{red\_cooperativa} permite que el
migrador resuelva \texttt{db-origen} y \texttt{db-destino} por nombre, sin
exponer los motores fuera del proyecto. El contenedor del migrador está
construido sobre \texttt{postgres:16} con \texttt{pgloader} 3.6.10,
\texttt{gettext-base} y \texttt{tzdata}; ya no requiere Python ni controladores
en el host.

\subsection{Permisos de archivos y secretos}

La clave privada TLS \texttt{server.key} tiene permisos \texttt{600} y
pertenece al usuario \texttt{postgres}. El certificado \texttt{server.crt}
tiene permisos \texttt{644}. El archivo \texttt{.env} contiene las
contraseñas de laboratorio, no se versiona y está incluido en
\texttt{.gitignore}. Los respaldos se guardan en \texttt{backups/}, también
excluido del control de versiones. En producción, las credenciales deben
gestionarse con un gestor de secretos y no con un archivo de texto.

\begin{table}[H]
  \centering
  \caption{Consideraciones del sistema operativo y del host}
  \label{tab:so}
  \contenidoConFuente{%
    \begin{tablaAPA}
      \begin{tabular}{@{}L{3.4cm}L{5.4cm}L{5.4cm}@{}}
        \toprule
        \textbf{Área} & \textbf{Configuración aplicada} & \textbf{Recomendación de producción} \\
        \midrule
        Aislamiento & Contenedores separados; red interna; puertos en \texttt{127.0.0.1}. & Mantener; añadir políticas de red y \emph{firewall}. \\
        Permisos & \texttt{server.key} en \texttt{600}; \texttt{.env} fuera de git. & Gestor de secretos; rotación; auditoría de accesos. \\
        Recursos & Sin límites explícitos en el laboratorio. & Definir \texttt{cpus}, \texttt{mem\_limit} y alertas. \\
        Registros & \texttt{docker logs}; \texttt{log\_line\_prefix}. & Rotación, recolección central y retención. \\
        Actualizaciones & Imágenes fijadas por versión mayor. & Escaneo de vulnerabilidades y parches programados. \\
        Zona horaria & \texttt{America/La\_Paz} en ambos motores. & Documentar y validar en reportes contables. \\
        Respaldos & Volúmenes nombrados; respaldo lógico cifrado. & Copia fuera del sitio y pruebas de restauración. \\
        \bottomrule
      \end{tabular}
    \end{tablaAPA}%
  }{Elaboración propia.}
\end{table}

\subsection{Zona horaria y consistencia temporal}

Ambos motores se configuraron con la zona \texttt{America/La\_Paz}. En MySQL,
la zona afecta la interpretación de \texttt{DATETIME} y las funciones de
fecha; en PostgreSQL, afecta la visualización de \texttt{TIMESTAMPTZ} y el
valor de \texttt{now()}. La migración asigna la zona en el ETL y el servidor
la conserva en las consultas. Para una entidad financiera, la consistencia
temporal es parte de la integridad: un pago no debe caer en el día anterior
o siguiente por un desajuste de zona.

\subsection{Actualizaciones y cadena de suministro}

Las imágenes se fijan por versión mayor (\texttt{mysql:5.7} y
\texttt{postgres:16}). En producción se recomienda fijar una versión exacta
(\emph{digest}) para evitar cambios inesperados, escanear las imágenes con una
herramienta de vulnerabilidades y actualizar con un proceso controlado. La
imagen de documentación (\texttt{texlive/texlive}) solo se usa en el perfil
\texttt{docs} y no forma parte del entorno operativo.

\section{Bitácora de problemas, decisiones y lecciones aprendidas}
\label{sec:bitacora}

\subsection{Problemas encontrados}

La Tabla~\ref{tab:bitacora} documenta los problemas reales que surgieron
durante la implementación. No se trata de una lista teórica: cada fila tiene
un síntoma, una causa, una solución y una evidencia. Documentar los errores es
parte del valor del proyecto, porque evita que el próximo equipo repita el
mismo camino.

\begingroup
\setstretch{1}\small\sloppy\emergencystretch=3em\relax
\begin{longtable}{@{}C{0.8cm}L{3.4cm}L{3.8cm}L{5.4cm}@{}}
  \caption{Bitácora de problemas y soluciones}\label{tab:bitacora}\\
  \toprule
  \textbf{\#} & \textbf{Síntoma} & \textbf{Causa} & \textbf{Solución aplicada} \\
  \midrule
  \endfirsthead
  \multicolumn{4}{c}{\textit{Continuación de la Tabla~\ref{tab:bitacora}}}\\
  \toprule
  \textbf{\#} & \textbf{Síntoma} & \textbf{Causa} & \textbf{Solución aplicada} \\
  \midrule
  \endhead
  \bottomrule
  \endlastfoot
  1 & El origen no tenía tablas al iniciar. & El volumen se creó antes de cargar el esquema; el entrypoint no se ejecutó. & Se recreó el volumen con \texttt{down -v} y se cargó el esquema en el primer arranque. \\
  2 & El seed inicial fallaba por cadenas sin comillas. & SQL generado a mano con valores de texto sin entrecomillar. & Se reemplazó por un generador Python determinista con \texttt{seed=42}. \\
  3 & \texttt{TRUNCATE ... RESTART IDENTITY} fallaba por permisos. & \texttt{admin\_dba} no era propietario de las secuencias de identidad. & Se eliminó \texttt{RESTART IDENTITY} y se usó \texttt{setval} al final. \\
  4 & PostgreSQL reiniciaba por certificado faltante. & La configuración apuntaba a una ruta no escribible o inexistente. & El entrypoint copia los certificados a \texttt{/var/lib/postgresql/ssl} y ajusta permisos. \\
  5 & El healthcheck generaba errores de \texttt{peer}. & Se ejecutaba como \texttt{root} con \texttt{-U postgres} por socket. & Se mantuvo el healthcheck; el contenedor queda \emph{healthy}; los errores son de autenticación \texttt{peer}, no de disponibilidad. \\
  6 & \texttt{auditor\_ro} podía ver datos sensibles. & La vista \texttt{vw\_cartera\_mora} exponía \texttt{numero\_identificacion}. & Se revocó el acceso a esa vista y se dejó solo \texttt{vw\_socios\_audit}. \\
  7 & El acceso remoto de \texttt{postgres} era posible. & \texttt{pg\_hba} incluía una regla \texttt{host} amplia. & Se eliminó la regla remota; solo \texttt{local peer}. \\
  8 & Las fechas perdían zona horaria. & \texttt{DATETIME} no almacena zona; el ETL no la asignaba. & Se usó \texttt{ZoneInfo("America/La\_Paz")} y \texttt{TIMESTAMPTZ}. \\
  9 & Los nuevos \texttt{INSERT} chocaban con IDs existentes. & Las secuencias no se reajustaron después de insertar IDs explícitos. & \texttt{setval} con \texttt{MAX(id)} en cada tabla. \\
  10 & El ETL no era idempotente. & Una segunda ejecución duplicaba filas. & \texttt{TRUNCATE} previo y snapshot consistente; verificación de conteos. \\
  11 & \texttt{postgres:16} es Debian 13 (trixie). & El paquete disponible es pgloader 3.6.10, no 3.6.9. & Se usó pgloader 3.6.10 (misma rama estable 3.6.x). \\
  12 & pgloader 3.6.10 verifica el certificado con \texttt{sslmode=require}. & \texttt{--no-ssl-cert-verification} no lo desactiva. & Se usó \texttt{SSL\_CERT\_FILE=/etc/postgresql/ssl/server.crt}. \\
  13 & pgloader no tiene opción \texttt{create no schema} para MySQL. & Crea un schema con el nombre de la base origen. & Se resolvió con \texttt{ALTER SCHEMA ... RENAME TO 'stg'}. \\
  14 & \texttt{ANALYZE} exige ser dueño en PostgreSQL 16. & \texttt{MAINTAIN} recién aparece en PG 17. & Se ejecuta como \texttt{postgres} por socket local. \\
  15 & pgloader escribe rechazos en \texttt{/tmp/pgloader/} por defecto. & Ruta por defecto no configurable sin parámetro. & Se usó \texttt{--root-dir /tmp/pgloader-work}. \\
  16 & Tags flotantes sin fijar por digest. & \texttt{mysql:5.7} y \texttt{postgres:16} pueden cambiar. & Mejora pendiente para producción: fijar por digest. \\
\end{longtable}
\fuenteAPA{Elaboración propia con base en la ejecución real y \texttt{evidence/}.}
\endgroup

\subsection{Decisiones de diseño}

La Tabla~\ref{tab:decisiones} resume las decisiones más importantes y su
justificación. Algunas fueron técnicamente debatibles; se documentan para que
el lector pueda evaluarlas y, si es necesario, cambiarlas con conocimiento de
causa.

\begin{table}[H]
  \centering
  \caption{Decisiones de diseño y su justificación}
  \label{tab:decisiones}
  \contenidoConFuente{%
    \begin{tablaAPA}
      \begin{tabular}{@{}L{3.8cm}L{4.4cm}L{6.0cm}@{}}
        \toprule
        \textbf{Decisión} & \textbf{Alternativa descartada} & \textbf{Justificación} \\
        \midrule
        ETL propio en Python & pgloader u otra herramienta genérica & Control explícito de tipos, zona horaria, identidades y verificación; el volumen es pequeño. \\
        Conservar nombres de origen & Renombrar columnas a un estándar nuevo & Reduce el impacto en la aplicación y facilita la comparación. \\
        \texttt{TIMESTAMPTZ} en \texttt{America/La\_Paz} & \texttt{TIMESTAMP} sin zona & Evita ambigüedad en pagos y reportes contables. \\
        \texttt{admin\_dba} sin \texttt{SUPERUSER} & Superusuario para todo & Limita el daño de un error o compromiso. \\
        Sin \texttt{RESTART IDENTITY} & \texttt{TRUNCATE ... RESTART IDENTITY} & Evita el error de propiedad de secuencias; \texttt{setval} es explícito. \\
        Puertos en \texttt{127.0.0.1} & Publicar en \texttt{0.0.0.0} & Reduce la exposición; el acceso remoto se hace por túnel. \\
        Vista enmascarada & Copia anonimizada completa & Permite auditoría sin duplicar datos ni romper referencias. \\
        \bottomrule
      \end{tabular}
    \end{tablaAPA}%
  }{Elaboración propia.}
\end{table}

\subsection{Lecciones aprendidas}

La primera lección es que la migración empieza por el diagnóstico, no por la
copia. Sin inventario, sin conteos y sin matriz de incompatibilidades, el ETL
es una caja negra. La segunda es que la seguridad debe diseñarse desde el
principio: si se deja para el final, se convierte en una capa de parches. La
tercera es que las pruebas deben ser automatizadas y comparar valores
esperados, no solo imprimir resultados. La cuarta es que la documentación de
los errores es tan valiosa como la de los aciertos: el problema de
\texttt{TRUNCATE ... RESTART IDENTITY} habría costado horas a otro equipo. La
quinta es que la reproducibilidad no es un lujo: es la única forma de afirmar
que el resultado es confiable.

\section{Lista de cotejo de seguridad}

La lista de cotejo de seguridad aplicada al destino se presenta en el
Anexo~C, sección C.2 (Tabla~\ref{tab:checklist}). Cada control tiene un estado
y una evidencia: los marcados como ``Cumple'' se verificaron en el entorno y
los marcados como ``Recomendado'' son mejoras para producción fuera del
alcance del laboratorio.

\section{Trazabilidad de evidencias}

La Tabla~\ref{tab:trazabilidad} relaciona cada punto exigido con la sección
del informe, el archivo de evidencia y el resultado. Esta tabla es la prueba
de que el informe no es una narración independiente de los datos: cada
afirmación tiene un origen verificable.

\begin{table}[H]
  \centering
  \caption{Trazabilidad de los puntos exigidos}
  \label{tab:trazabilidad}
  \contenidoConFuente{%
    \begin{tablaAPA}
      \begin{tabular}{@{}C{0.7cm}L{4.6cm}L{3.8cm}L{4.6cm}@{}}
        \toprule
        \textbf{\#} & \textbf{Punto exigido} & \textbf{Evidencia} & \textbf{Resultado} \\
        \midrule
        1 & Portada, resumen, objetivos y alcance. & Este documento & Secciones 1.1--1.7. \\
        2 & Cliente, problema, solución y ODS. & Este documento & Secciones 1.2--1.5. \\
        3 & Diagnóstico del origen. & \texttt{evidence/01} & 8 tablas, 2 vistas, 0.28 MB. \\
        4 & Hardening y verificación. & \texttt{evidence/03} & 75 PASS, 0 FAIL. \\
        5 & Plan de migración y contingencia. & Este documento & Sección 3.1. \\
        6 & Ejecución y salidas reales. & \texttt{evidence/02} & 1\,469 filas, 0 diferencias. \\
        7 & Pruebas en destino. & \texttt{evidence/03} & Conteos, FK y rendimiento. \\
        8 & Confidencialidad y escalabilidad. & Este documento & Sección 3.8. \\
        9 & Consideraciones del SO. & Este documento & Sección 3.9. \\
        10 & Bitácora y lecciones. & Este documento & Sección 3.10. \\
        11 & Conclusiones, recomendaciones y bibliografía. & Este documento & Capítulo 5 y referencias. \\
        12 & Anexos y lista de cotejo. & Anexos A--C & Scripts completos. \\
        13 & Guion de demostración. & \texttt{GUION\_DEMO\_EN\_VIVO.md} & 15--20 min. \\
        \bottomrule
      \end{tabular}
    \end{tablaAPA}%
  }{Elaboración propia.}
\end{table}

\section{Reproducibilidad y transferencia}

\subsection{Procedimiento de ejecución}

El procedimiento completo se ejecuta en cuatro pasos. Primero, se copia el
archivo \texttt{.env.example} a \texttt{.env} y se ajustan las credenciales.
Segundo, se levanta el entorno con \texttt{docker compose up -d --build}.
Tercero, se ejecutan los scripts en orden. Cuarto, se revisan las evidencias y
se compila el informe. La secuencia exacta se muestra a continuación.

\begin{lstlisting}[style=terminal,caption={Secuencia de reproducción del proyecto},label={lst:reproduccion}]
cp .env.example .env
docker compose up -d --build
./scripts/01_audit_origin.sh
./scripts/02_migrate_and_backup.sh
./scripts/03_verify_all.sh
docker compose --profile docs run --rm docs make
\end{lstlisting}

\subsection{Resultados esperados}

Si el entorno se levanta desde cero, los resultados esperados son: 8 tablas y
2 vistas en el origen; 55 socios, 66 cuentas, 50 solicitudes, 32 préstamos,
516 cuotas y 740 transacciones; 1\,469 filas migradas; 75 pruebas en
\texttt{PASS}; y un PDF en \texttt{docs/main.pdf}. Cualquier diferencia debe
investigarse antes de aceptar la migración. La semilla fija del generador
garantiza que los datos sean los mismos en cada ejecución.

\subsection{Transferencia}

El proyecto se entrega con la siguiente estructura: \texttt{src/} (esquemas,
seed, ETL y configuración), \texttt{scripts/} (automatización),
\texttt{evidence/} (salidas reales), \texttt{backups/} (respaldos cifrados),
\texttt{docs/} (informe, anexos y guion) y \texttt{docker-compose.yml}. Las
credenciales de laboratorio están en \texttt{.env.example}; las reales no se
versionan. El lector puede reproducir el entorno en cualquier host Linux con
Docker y Docker Compose, sin instalar MySQL ni PostgreSQL en el sistema.
